\documentclass[11pt,a4paper]{article}

% ─── Packages ────────────────────────────────────────────────────────────────
\usepackage[utf8]{inputenc}
\usepackage[T1]{fontenc}
\usepackage{geometry}
\geometry{a4paper, margin=1in}
\usepackage{booktabs}
\usepackage{tabularx}
\usepackage{xcolor}
\usepackage{hyperref}
\usepackage{enumitem}
\usepackage{tcolorbox}
\usepackage{titlesec}
\usepackage{amssymb}

% ─── Colours & Styling ────────────────────────────────────────────────────────
\definecolor{primaryblue}{RGB}{0, 102, 204}
\definecolor{systembg}{RGB}{240, 248, 255}
\definecolor{systemframe}{RGB}{176, 196, 222}
\definecolor{riskbg}{RGB}{255, 245, 238}
\definecolor{riskframe}{RGB}{210, 180, 140}
\definecolor{warningbg}{RGB}{255, 243, 205}
\definecolor{warningframe}{RGB}{255, 193, 7}
\definecolor{checkbg}{RGB}{232, 245, 233}
\definecolor{checkframe}{RGB}{76, 175, 80}
\definecolor{comparativebg}{RGB}{243, 229, 245}
\definecolor{comparativeframe}{RGB}{156, 39, 176}

\hypersetup{
    colorlinks=true,
    linkcolor=primaryblue,
    filecolor=magenta,
    urlcolor=primaryblue,
}

% ─── Tcolorbox environments ────────────────────────────────────────────────────
\tcbset{
    systembox/.style={
        detach title,
        breakable,
        colback=systembg,
        colframe=systemframe,
        left=12pt, right=12pt, top=10pt, bottom=10pt,
        title={#1},
        fonttitle=\bfseries\large,
        title style={colback=systemframe!40!white},
        arc=4pt,
    },
    riskbox/.style={
        detach title,
        breakable,
        colback=riskbg,
        colframe=riskframe,
        left=12pt, right=12pt, top=10pt, bottom=10pt,
        title={#1},
        fonttitle=\bfseries,
        title style={colback=riskframe!40!white},
        arc=4pt,
    },
    methodbox/.style={
        detach title,
        breakable,
        colback=checkbg,
        colframe=checkframe,
        left=12pt, right=12pt, top=10pt, bottom=10pt,
        title={#1},
        fonttitle=\bfseries,
        title style={colback=checkframe!40!white},
        arc=4pt,
    },
    warnbox/.style={
        detach title,
        breakable,
        colback=warningbg,
        colframe=warningframe,
        left=12pt, right=12pt, top=10pt, bottom=10pt,
        title={#1},
        fonttitle=\bfseries,
        title style={colback=warningframe!40!white},
        arc=4pt,
    },
    comparativebox/.style={
        detach title,
        breakable,
        colback=comparativebg,
        colframe=comparativeframe,
        left=12pt, right=12pt, top=10pt, bottom=10pt,
        title={#1},
        fonttitle=\bfseries,
        title style={colback=comparativeframe!40!white},
        arc=4pt,
    },
}

% ─── Document ─────────────────────────────────────────────────────────────────
\begin{document}

% ─── COVER ────────────────────────────────────────────────────────────────────
\begin{titlepage}
\vspace*{2cm}
\begin{center}
{\LARGE\bfseries Data Scientist \& AI Engineer}\\[0.4em]
{\large\itshape Portfolio Build Plan}\\[1em]
{\LARGE\bfseries Two-System Master Portfolio}\\[0.6em]
{\large GenAI Engineering \textbullet\ Production Systems \textbullet\ Research Rigor}\\[2cm]
\end{center}

\noindent\textbf{Core Strategy:} Depth over breadth. Two polished, production-grade systems built end-to-end, covering the full modern GenAI stack — from ingestion, retrieval, and secure agentic orchestration to alignment, rigorous evaluation, and observability. 

\vspace{1.5cm}
\noindent\textbf{Why cut Projects 3, 4 and SSL?}
\begin{itemize}
    \item\textbf{Project 3 (BI/SQL):} Already covered by your Master in BI, DP-600 cert, and professional experience. Adds no new market signal.
    \item\textbf{Project 4 (n8n job scraper):} Too cliché; screams \emph{job seeker} rather than \emph{AI system architect}.
    \item\textbf{SSL project (DINOv2):} Your CV already has MoCo v2 and BYOL with real-world metrics. A hobby demo would dilute — not reinforce — that story.
\end{itemize}

\vspace{1.5cm}
\begin{center}
\small
\begin{tabularx}{0.9\textwidth}{l X}
    \toprule
    \textbf{System 1} & Railway Regulatory Intelligence Platform (Multimodal RAG, Agents, MCP, Security, LLMOps) \\
    \textbf{System 2} & Small-LLM Alignment (GRPO vs. DPO Comparative Analysis, Training MLOps) \\
    \textbf{Corpus}   & SNCF regulatory documents (DILA legifrance, EU directives) \\
    \textbf{Timeline} & Sequential — finish System 1 before System 2 \\
    \textbf{Goal}     & Professional portfolio, LinkedIn/YouTube showcase, CV signal \\
    \bottomrule
\end{tabularx}
\end{center}
\end{titlepage}

\newpage

% ─── STABILITY FILTER ──────────────────────────────────────────────────────────
\section*{Stability Filter Applied — Final Tech Stack}

\vspace{0.5em}
\noindent
\begin{tabularx}{\textwidth}{@{} l l X @{}}
\toprule
\textbf{Component} & \textbf{Choice} & \textbf{Rationale} \\
\midrule
Orchestration & LangGraph & Stateful, explainable graphs; supports human-in-the-loop \\
Observability & Langfuse & Full LLM tracing; production-ready open-source \\
LLM Layer     & OpenAI GPT-4o-mini (API) & Cost-effective, reliable quality \\
Embedding     & text-embedding-3-small   & Cheap, good quality, OpenAI native \\
Vector store  & Qdrant (local Docker)    & Lightweight, fast, free, GPU-accelerated \\
Agent runtime & LangGraph native         & Full control over state, recovery, and tools \\
MCP server    & FastMCP / official SDK & Secure standard protocol; real tools vs. mocks \\
Fine-tuning   & Unsloth + TRL            & 4-bit LoRA; GRPO and DPO viability on 8GB GPU \\
Eval          & RAGAS + Giskard          & Industry-standard; statistical eval frameworks \\
UI            & Streamlit (shared)       & Fast prototyping; deployable to Streamlit Cloud \\
Packaging     & Docker + GitHub Actions  & Non-negotiable enterprise standard \\
\bottomrule
\end{tabularx}
\vspace{1em}

% ─── SYSTEM 1 ──────────────────────────────────────────────────────────────────
\newpage
\section*{SYSTEM 1: Railway Regulatory Intelligence Platform}

\subsection*{Why This System Exists}
SNCF operates under a dense, continuously evolving regulatory environment. Customer-facing and internal teams need reliable, cited answers drawn directly from official regulatory text — complete with deterministic validation, failure recovery, and secure tool usage. 

\textbf{Job market signal:} Hybrid RAG, Agentic orchestration, human-in-the-loop approval, state recovery, observability, and enterprise security.

% ── W1: Ingestion ──
\subsection*{Week 1 — Data Ingestion \& Retrieval Engineering}

\begin{tcolorbox}[systembox={ Step-by-Step}]
\begin{itemize}[leftmargin=*]
    \item \textbf{D1–2: Dataset construction}\\
          Download 50–200 PDFs from DILA legifrance (Code des transports, EU rail directives). Deduplicate, strip noise, resolve encoding issues.
    \item \textbf{D3: Chunking strategies (Ablation Study)}\\
          Compare recursive character split vs. semantic chunking vs. structure-aware chunking (section-header aware).
  \item \textbf{D4–5: Hybrid Search \& Reranking (Qdrant)}\\
      Embed with \texttt{text-embedding-3-small}. Implement BM25 sparse vectors + dense vectors with Reciprocal Rank Fusion (RRF). Pass the top 20 results through a cross-encoder (\texttt{BAAI/bge-reranker-v2-m3}) to output the final top-5 most relevant chunks. 
\end{itemize}
\end{tcolorbox}

% ── W2: Eval + Observe + Ship ──
\subsection*{Week 2 — Evaluation, Observability \& UI}

\begin{tcolorbox}[systembox={ Step-by-Step}]
\begin{itemize}[leftmargin=*]
    \item \textbf{D1–2: Test set + RAGAS evaluation}\\
          Curate a golden Q\&A set. Run RAGAS (Faithfulness, Answer Relevancy, Context Precision). Log baseline scores.
    \item \textbf{D3: Langfuse integration}\\
          Instrument all calls. Trace retrieval $\rightarrow$ generation chains.
    \item \textbf{D4–7: Enterprise UI (Open WebUI)}\\
          Instead of a basic Streamlit app, deploy \textbf{Open WebUI} via Docker Compose alongside your backend. Expose your LangGraph agent as an API to Open WebUI to demonstrate a production-ready frontend with native citation rendering.
\end{itemize}
\end{tcolorbox}

% ── W3: Agentic ──
\subsection*{Week 3 — Enterprise Agentic Workflow (LangGraph)}

\begin{tcolorbox}[systembox={ Step-by-Step: State, Recovery \& Approval}]
\begin{itemize}[leftmargin=*]
    \item \textbf{D1: Graph Architecture}\\
          Map conditional routing: User Query $\rightarrow$ Router $\rightarrow$ [Retriever | Summarizer | Escalation].
    \item \textbf{D2–4: State Persistence \& Failure Recovery}\\
          Implement \texttt{MemorySaver} (SQLite checkpointer) to persist conversation state. Add a \texttt{max\_retries} loop for hallucinations (deterministic validation check $\rightarrow$ retry if failed).
    \item \textbf{D5–7: Human-in-the-Loop (HITL)}\\
          Add an \texttt{interrupt\_before} node for sensitive actions. The graph pauses, saves state, and waits for human approval via the UI before proceeding.
\end{itemize}
\end{tcolorbox}

% ── W4: MCP + Action Layer (n8n) ──
\subsection*{Week 4 — MCP Security \& n8n Automation}

\begin{tcolorbox}[systembox={ Step-by-Step: Secure Tools \& Workflows}]
\begin{itemize}[leftmargin=*]
    \item \textbf{D1–2: Secure MCP Server}\\
          Expose a live API tool (e.g., fetching live SNCF traffic alerts). \textbf{Security focus:} Implement strict environment variable isolation and scope-limited read-only access for the agent.
    \item \textbf{D3–4: n8n Workflow Automation Integration}\\
          We can ignore using n8n completely and use directly LangGraph, it depends after i learn how to apply them and how they work: Connect the LangGraph HITL output to an \textbf{n8n} webhook. Once a human approves the generated regulatory report, the graph triggers n8n to automatically format the output and route it to a mock Slack channel or email sandbox.
    \item \textbf{D5–7: DeepEval \& Final Polish}\\
          Run agent trajectory evaluations (Did it use the MCP tool correctly? Did it wait for human approval before triggering n8n?). Finalize the Docker Compose network and ship.
\end{itemize}
\end{tcolorbox}

% ── EXTENSION 1: Multimodal Vision ──
\subsection*{Phase 2 Extension — Multimodal Regulatory Understanding}
\textit{Replaces the need for a separate Computer Vision project by retrofitting System 1 to handle visual data, tying perfectly into your existing CV background.}

\begin{tcolorbox}[systembox={ Step-by-Step: Vision Integration}]
\begin{itemize}[leftmargin=*]
    \item \textbf{D1–2: Visual Corpus Curation}\\
          Isolate regulatory PDFs containing complex visual constraints (railway safety pictograms, complex DILA routing tables, scanned annexes).
    \item \textbf{D3–4: Vision-Native Ingestion (ColPali)}\\
          Implement \texttt{ColPali} to process PDF pages directly as image patches rather than mangling them through text extractors like PyPDF. Store image embeddings in a dedicated vector index.
    \item \textbf{D5–7: VLM Agent Integration}\\
          Add a routing node in LangGraph to detect visual queries. Pass retrieved image patches to a Vision-Language Model (e.g., GPT-4o or Claude 3.5 Sonnet) to answer questions based directly on the visual context of the regulatory documents.
\end{itemize}
\end{tcolorbox}

% ── EXTENSION 2: MLOps Layer ──
\subsection*{Phase 3 Extension — Cross-System MLOps \& CI/CD Backbone}
\textit{Applied retroactively to elevate this from a "notebook demo" to production-grade engineering.}

\begin{tcolorbox}[systembox={ Step-by-Step: CI/CD \& Drift Monitoring}]
\begin{itemize}[leftmargin=*]
    \item \textbf{D1–2: GitHub Actions CI/CD Pipeline}\\
          Write workflow files (\texttt{.github/workflows/main.yml}) that build your Docker containers on every pull request to \texttt{main}.
    \item \textbf{D3–4: Automated Regression Testing}\\
          Integrate RAGAS/DeepEval into the GitHub Actions runner. On every PR, the pipeline runs 20 golden queries against the proposed code. If Faithfulness or Context Precision drops below 0.85, the pipeline automatically fails the PR merge.
    \item \textbf{D5–7: Production Drift Monitoring}\\
          Set up a continuous monitor to detect semantic data drift. If users start asking about new regulatory domains not present in the Qdrant index, log alerts to a mock Slack channel (re-using your n8n setup) to notify the data team to update the ingestion pipeline.
\end{itemize}
\end{tcolorbox}


% ─── SYSTEM 2 ─────────────────────────────────────────────────────────────────
\newpage
\section*{SYSTEM 2: Small-LLM Alignment (GRPO vs. DPO)}

\subsection*{Why This System Exists}
On-device or cloud-small LLMs (1.5B - 3B) are economically superior for production, but require alignment. Demonstrating modern RLHF (GRPO) alongside a contrastive baseline (DPO) proves you understand the mathematical tradeoffs of model alignment, not just how to call an API. Applying strict Training MLOps proves you can manage and deploy these alignments in an enterprise setting.

% ── Hardware ──
\begin{tcolorbox}[warnbox={️ Hardware Constraints}]
\begin{itemize}[leftmargin=*]
    \item \textbf{GPU:} 8GB VRAM. Use 4-bit QLoRA via Unsloth.
    \item \textbf{RAM:} 32 GB system RAM recommended.
    \item \textbf{Compute Strategy:} Train adapters only. Keep batch sizes small (group size 4 for GRPO).
\end{itemize}
\end{tcolorbox}

% ── Reward Modeling ──
\subsection*{Reward Model Engineering \& Training MLOps}

\begin{tcolorbox}[systembox={ Step-by-Step}]
\begin{itemize}[leftmargin=*]
    \item \textbf{D1–2: Reward function design (Rule-based + Judge)}\\
          1) Accuracy (correct clause citation), 2) Format compliance, 3) Length penalty.
    \item \textbf{D3: Reward Model Validation}\\
          Ensure the reward correlates with human preference on a golden set before training begins to prevent reward hacking.
    \item \textbf{D4–5: Training Infrastructure (MLOps)}\\
          Implement robust training MLOps using \textbf{Weights \& Biases (W\&B)} for tracking reward metrics, loss curves, and GPU utilization during the RL phase. Configure automated model registry versioning and canary deployment scripts.
\end{itemize}
\end{tcolorbox}

% ── Comparative Alignment ──
\subsection*{The Professional Differentiator: GRPO vs. DPO Comparison}

\begin{tcolorbox}[comparativebox={ Optional Comparative Analysis: GRPO vs. DPO}]
To exhibit true senior research rigor, optionally train two separate LoRA adapters on your preference dataset and compare the methodologies:

\begin{itemize}[leftmargin=*]
    \item \textbf{Pipeline A: DPO (Direct Preference Optimization)}\\
          Offline contrastive learning. Feed the model chosen/rejected pairs.
          \textit{Expectation:} Faster to train, mathematically stable, good for formatting, but lacks reasoning depth.
    \item \textbf{Pipeline B: GRPO (Group Relative Policy Optimization)}\\
          Active RL. The model generates 4-8 answers, scores them via the reward function, and learns via self-exploration (no value network needed).
          \textit{Expectation:} Harder to stabilize, uses slightly more memory during rollout, but drastically improves logical reasoning and adherence to complex constraints.
    \item \textbf{The Deliverable:} A comparison table in your README showing Training Time, Memory Usage, Final Reward Score, and qualitative formatting accuracy between the Base Model, DPO, and GRPO.
\end{itemize}
\end{tcolorbox}

% ── Regression & Stats ──
\subsection*{Methodological Rigor}

\begin{tcolorbox}[methodbox={✅ Evaluation \& Statistical Honesty}]
\begin{itemize}[leftmargin=*]
    \item \textbf{Regression Testing:} Run the pre-trained base model on standard prompts before and after training to ensure you haven't caused "capability collapse" (forgetting how to speak normally).
    \item \textbf{Statistical Confidence:} On the held-out test set, report 95\% Confidence Intervals via bootstrap resampling. Never report just a mean score.
    \item \textbf{Limitations Section:} Explicitly state the boundaries (dataset size, lack of human preference data, domain scope).
\end{itemize}
\end{tcolorbox}

% ── MLOps & Model Serving ──
\subsection*{Production MLOps \& Model Serving Integration}

\begin{tcolorbox}[systembox={ Step-by-Step: Registry \& Inference Endpoint}]
\begin{itemize}[leftmargin=*]
    \item \textbf{D1–2: Model Registry \& Versioning}\\
          Use \textbf{Weights \& Biases (W\&B) Artifacts} or Hugging Face Hub to strictly version your LoRA adapters. Tie each model version to its specific git commit and dataset hash.
    \item \textbf{D3–4: Automated Evaluation Pipeline}\\
          Build a pipeline script that automatically pulls the latest adapter, merges it with the base model, and runs your regression tests against the held-out dataset, logging the final metrics back to W\&B.
    \item \textbf{D5–7: Production Serving (vLLM / Ollama)}\\
          Package the aligned model into an \textbf{Ollama} Modelfile or deploy it via \textbf{vLLM} in a Docker container, proving you know how to transition a fine-tuned adapter from a training notebook to a high-throughput inference endpoint.
\end{itemize}
\end{tcolorbox}

% ─── CROSS-SYSTEM OBSERVABILITY ────────────────────────────────────────────────
\newpage
\section*{Cross-System Observability Layer}

\begin{tcolorbox}[systembox={ Langfuse Observability — Applied to Both Systems}]
\begin{tabularx}{\textwidth}{@{} l X @{}}
\toprule
\textbf{System 1 (RAG/Agents)} & Log: Query, top-k chunks, state transitions, MCP tool calls, agent latency, and human-approval timestamps. \\
\midrule
\textbf{System 2 (Alignment)} & Log: Per-step reward, KL penalty, gradient norm, evaluation test samples, and prediction traces. \\
\bottomrule
\end{tabularx}
\vspace{1em}
\textit{Both systems are instrumented using the same self-hosted Langfuse project, proving cross-platform architectural consistency.}
\end{tcolorbox}

% ─── POSTING STRATEGY ──────────────────────────────────────────────────────────
\section*{LinkedIn / YouTube Posting Strategy}

\bigskip
\noindent
\begin{tabularx}{\textwidth}{@{} l l X @{}}
\toprule
\textbf{Post} & \textbf{Platform} & \textbf{Content} \\
\midrule
1 & LinkedIn & Corpus, Chunking \& Multimodal Vision: Why naive RAG fails on regulations. \\
2 & LinkedIn & LangGraph Architecture: State recovery and Human-in-the-Loop. \\
3 & YouTube & System 1 Demo: Secure MCP integration and cited answers. \\
4 & LinkedIn & LLMOps \& CI/CD: Automated RAGAS testing on GitHub Actions. \\
5 & LinkedIn & GRPO vs DPO: Why I chose active RL over offline preference matching. \\
6 & LinkedIn & Training MLOps \& Statistical Honesty: W\&B, effect sizes, and avoiding reward hacking. \\
\bottomrule
\end{tabularx}

\vspace{2cm}
\begin{center}
    \textit{"Most portfolios show what tools a developer can import. \\ This portfolio shows how an engineer builds, secures, aligns, and evaluates systems in production."}
\end{center}
\end{document}

\end{document}